
@@M0@@
\uline{(où} @@M1@@\uline{) est une application polynôme. Si} @@M2@@\uline{, alors} @@M3@@ \uline{est de degré} @@M4@@.

\emph{Démonstration} : Soit @@M5@@ une suite d'éléments de @@M6@@ avec @@M7@@ (@@M8@@). Alors la formule binomiale donne pour tout @@M9@@ :
@@M10@@
où la somme à droite est finie puisque les éléments @@M11@@ sont nilpotents. Mais les coefficients @@M12@@ sont des polynômes à coefficients rationnels dans les variables @@M13@@ de degré @@M14@@, de sorte que @@M15@@ est bien une application polynôme. Si @@M16@@, alors @@M17@@ pour @@M18@@. Donc @@M19@@ est de degré @@M20@@, C.Q.F.D.

\subsection*{2.2.}
Considérons maintenant un schéma noethérien @@M21@@ et un élément @@M22@@ de @@M23@@ (1.1.1). En associant à tout @@M24@@ l'élément @@M25@@ de @@M26@@ on obtient un homomorphisme de @@M27@@-modules
@@M28@@
qui se factorise, compte tenu de (1.3.2), par l'épimorphisme canonique
@@M29@@
